% from ega4-2-fr.tex
\oldpage[IV-2]{54}

raisonnement que dans (4.1.2.1) montre que si $\dim(U_i)=n_i$, il existe un
$k$-morphisme $g_i:U_i\to\mathbf V_k^{n_i}$ qui est fini et dominant, donc
surjectif (II, 6.1.10). Soit $n=\dim(X)=\sup_i n_i$~; pour chaque $n_i$ il y a
donc un morphisme $h_i:\mathbf V_k^{n_i}\to\mathbf V_k^n$ qui est une
immersion fermée correspondant à l'homomorphisme canonique
$k[T_1,\ldots,T_n]\to k[T_1,\ldots,T_{n_i}]$, et qui est l'identité pour
$n_i=n$. Comme $U$ est somme des $U_i$, on prendra pour $g$ le morphisme qui
dans chaque $U_i$ coïncide avec $h_i\circ g_i$~; il est évidemment fini et
surjectif. Lorsque $n'\geqslant n$, on aura un morphisme fini
$U\to\mathbf V_k^{n'}$ en composant l'immersion fermée canonique
$\mathbf V_k^n\to\mathbf V_k^{n'}$ avec $g$~; lorsque $n'\leqslant n$, on
compose de même avec $g$ le morphisme canonique
$p:\mathbf V_k^n\to\mathbf V_k^{n'}$ correspondant à l'injection canonique
$k[T_1,\ldots,T_{n'}]\to k[T_1,\ldots,T_n]$, en notant que $p$ est fidèlement
plat, donc \emph{surjectif}.
\end{proof}

\begin{remark}[4.1.3]
\label{IV.2.4.1.3-fr}
Le corollaire (4.1.2.1) montre que dans la formule (4.1.1.1), on peut supposer
que $x$ parcourt l'ensemble de \emph{tous les points} de $X$.
\end{remark}

\begin{corollary}[4.1.4]
\label{IV.2.4.1.4-fr}
Soient $X$ un préschéma localement de type fini sur un corps $k$, $K$ une
extension de $k$~; alors $\dim(X\otimes_k K)=\dim(X)$.
\end{corollary}

\begin{proof}
On peut évidemment se borner au cas où $X$ est de type fini sur $k$~; alors le
morphisme $u:X\otimes_k K\to X$ est fidèlement plat (2.2.13, (i)), donc si
$U$ est un ouvert partout dense dans $X$, $u^{-1}(U)=U\otimes_k K$ est dense
dans $X\otimes_k K$ (2.3.10)~; si $g:U\to\mathbf V_k^n$ est fini et
surjectif, il en est de même de
$g_{(K)}:U\otimes_k K\to\mathbf V_K^n$ (I, 3.5.2 et II, 6.1.5), d'où le
corollaire.
\end{proof}

\begin{corollary}[4.1.5]
\label{IV.2.4.1.5-fr}
Soient $X$ et $Y$ deux préschémas localement de type fini sur un corps $k$~;
alors $\dim(X\times_k Y)=\dim(X)+\dim(Y)$.
\end{corollary}

\begin{proof}
Il suffit de démontrer que si $U$ (resp. $V$) est un voisinage affine d'un
point $x$ (resp. $y$) de $X$ (resp. $Y$) tel que
$\operatorname{deg.tr.}_k k(x)=m=\dim U$ et
$\operatorname{deg.tr.}_k k(y)=n=\dim V$, alors on a
$\dim(U\times_k V)=m+n$, autrement dit, on peut supposer en vertu de (4.1.2)
qu'il existe des $k$-morphismes finis surjectifs
$f:X\to\mathbf V_k^m$, $g:Y\to\mathbf V_k^n$~; alors
$f\times g:X\times_k Y\to\mathbf V_k^{m+n}$ est fini et surjectif
(I, 3.5.2 et II, 6.1.5), d'où le corollaire.
\end{proof}

\subsection{Cycles premiers associés sur les préschémas algébriques}
\label{subsection:IV.2.4.2-fr}

\begin{proposition}[4.2.1]
\EGAReaderTarget{ega4.c04.prop.4-2-1}{IV.2.4.2.1-fr}
Soient $K$ et $L$ deux extensions d'un corps $k$, telles que $K\otimes_k L$
soit noethérien. Alors les idéaux premiers associés à $K\otimes_k L$ sont
minimaux, et si $E$ est le corps résiduel de l'anneau local d'un tel idéal, on
a
\begin{equation}
\tag{4.2.1.1}
\label{IV.2.4.2.1.1-fr}
\operatorname{deg.tr.}_K E=\operatorname{deg.tr.}_k L,
\qquad
\operatorname{deg.tr.}_L E=\operatorname{deg.tr.}_k K,
\end{equation}
d'où
\begin{equation}
\tag{4.2.1.2}
\label{IV.2.4.2.1.2-fr}
\operatorname{deg.tr.}_k E
=\operatorname{deg.tr.}_k K+\operatorname{deg.tr.}_k L.
\end{equation}
\end{proposition}

\begin{proof}
On sait que $K$ est une extension algébrique d'une extension transcendante
pure $K'=k(\mathbf t)$, où $\mathbf t=(t_\alpha)_{\alpha\in I}$ est une famille
d'indéterminées~; l'anneau $k[\mathbf t]\otimes_k L=L[\mathbf t]$ est intègre,
donc il en est de même de $K'\otimes_k L$, qui en est un anneau de fractions,
et le corps des fractions de $K'\otimes_k L$ est $L(\mathbf t)$~; on a le
diagramme commutatif d'homomorphismes canoniques

